\begin{abstract}
Scientific ideas are often represented by embedding the papers that express them. These embeddings help find similar papers, but once an idea is reduced to a vector, the vector itself cannot say what idea it represents without external evidence such as nearby papers or a separately trained decoder. We propose to use the LoRA adapters generated by hypernetworks as document embeddings that represent scientific idea in papers. The adapters can be loaded into a frozen language model and queried in natural language. As a result, the original papers as well as averages, interpolations, and other constructed points in the space can be directly decoded and interpreted. 
For example, decoding the average of papers in a physics subfield produces a label close to the official field name, and midpoints between two papers decode to descriptions that combine both sources.
We evaluated whether LoRA embeddings can be used as representation for similarity search and clustering in the tasks of next-paper prediction, topic classification, collaboration prediction, and author-name disambiguation, finding that they are on par with SPECTER2, Instructor, EmbeddingGemma, and GTE without sacrificing the decodability. 
Our results show that these embeddings allow not only retrieval but also direct and arithmetic interpretability and operability of scientific ideas.
\end{abstract}
